\section{Conclusion}
In this work, we extended the Baleen (FAST~'24) machine learning–guided flash caching framework with a focus on improving eviction behavior under dynamic backend load conditions. While Baleen demonstrated strong results in admission and prefetching, we still found that its eviction logic had limited adaptability when workloads shift or contention increased. To address this gap, we introduced two enhanced eviction mechanisms: DT-SLRU and Episode-Deadline Eviction (EDE). DT-SLRU adjusts promotion decisions based on disk service time, while EDE prioritizes blocks according to estimated remaining usefulness and protects high-impact blocks from premature eviction. Through sensitivity analysis and ablation studies on key parameters, we observed that proper tuning of the DT promotion threshold ($\tau_{DT}$) and protected-cap ratio enables significant performance gains. In particular, our experiments show that setting $\tau_{DT} \approx 2$ and using a protected-cap of approximately 0.5 can reduce Peak Disk-head Time by up to \textbf{45\%} compared to the baseline configuration, indicating that eviction policy behavior is highly influential in determining overall system stability and responsiveness.

Although our enhancements improve adaptability, there are still open opportunities for future work. One direction is to integrate online or workload-aware parameter tuning rather than relying on static thresholds. Another avenue is to co-train eviction and admission models end-to-end, allowing the system to respond cohesively to changing load characteristics. Additionally, deploying these techniques in real storage environments would provide deeper insights into long-term flash wear, operational cost, and performance trade-offs. Overall, our findings highlight that combining ML-guided admission with adaptive eviction policies can create a more robust and cost-effective flash caching system that can maintain strong performance even under unpredictable or bursty workloads.
